% ch4_b §4.3–4.5 (书页 129–138 / p143–p152)

\section{内聚性}

固体理论的一个主要目标，是说明固体何以是其所是——也就是说，要证明把一大堆，比如说，Na 原子放进一个盒子里，它们会凝聚成一块金属钠晶体。这个目标过于雄心勃勃；但我们还是可以研究固体能量学的某些侧面。特别是，我们可以尝试计算固体的内聚能（cohesive energy）——从分离原子集合体的能级量起的能量——并将它与实验相比较，或者与由同样原子构成的其他假想固体（点阵常数不同，或晶体晶格不同）的能量相比较。

例如，注意到分子晶体的内聚能很小——量级为每分子 0·1 eV——我们就证实了 van der Waals 力之弱。氢键型晶体（典型如 H$_2$O）的内聚要高一些——比如每分子 0·5 eV。这两类固体的熔点都很低。金属则更够得上典型的“固体”，内聚能介于每原子 1 至 5 eV 之间；而离子型和共价型固体结合得非常牢固，能量达每原子 10 eV 的量级。

着手计算内聚能所依循的原理，从前几章已可看出。例如对离子晶体，我们先为该结构算出 Madelung 常数（式 (2.27)），从而把静电力的贡献确定为
\begin{equation*}
\mathcal{E}=-\frac{1}{2}N\alpha\frac{e^{2}}{a}.
\tag{4.1}
\end{equation*}
这个能量显然趋于使固体塌缩，即使点阵常数 $a$ 减小。但这个力受到离子实之间排斥的对抗——这种排斥常用如下形式的势来代表：
\begin{equation*}
\phi(R)=B\exp(-R/\rho)
\tag{4.2}
\end{equation*}
其中 $R$ 为离子间的距离。

把这两项拼在一起、把固体的总能量估计为 $a$ 的函数，是一个不足道的计算。它可以同总内聚能以及固体压缩系数的实验值相比较。遗憾的是，对给定的一对离子，参数 $B$ 与 $\rho$ 不可能从第一性原理以任何高精度算出；我们至多只能导出一组经验值，它们或多或少地与对各种不同晶体——其中离子按一切可能的组合配成对——的观测相一致。这个唯象方案于是可以解释为给出若干元素的离子半径（ionic radii）。

共价固体的内聚能是一个复杂得多的问题。原则上，可以通过计算能带结构、再求出每个电子态能量的贡献而把它找出来。实际上，更容易的做法是证实：晶体的总能量相当接近各共价键能量之和，而晶格谱可以由在声子模中弯曲和/或拉伸这些键所需的力导出。就金刚石而言，这种符合是极好的。如果我们懂得了共价键的能量学，那么在理解半导体内聚能的道路上，我们就已经走了很远了。

既然共价固体的内聚几乎完全依赖于键的化学饱和性，结构就不一定非要具有晶体序不可。在一种没有长程序的网络中，常常也能让每个原子以非常接近正确的几何组态获得正确数目的近邻。例如，普通的硅石玻璃（silica glass）就是这样的集合体：每个 Si 原子都位于一个由氧构成的四面体的中心，而每个氧原子几乎正好位于两个 Si 原子的连线上（图 77）。这个网络在化学上和动力学上都显然是稳定的，但应当如何从能带的观点描述它的内聚，却不清楚。

%FIG{77}: 共价键网络：(a) 晶态；(b) 玻璃态。这是 SiO$_2$ 的二维类似物。

金属内聚（metallic cohesion）则是一个微妙得多的现象。大体上我们可以说，把每个电子从它在特定原子上的定域中解放出来，就增大了其位置坐标的不确定度；于是按 Heisenberg 原理，动量的散布便减小。这降低了电子的平均动能，从而把系统结合在一起。但是，总内聚能的计算属于整个金属理论中最难做对、也最难在原则上验证的问题之列。为了计入电子气体中的关联与交换这类复杂的多体效应（见 §5.8），切分这块“蛋糕”的办法有很多种。

在最简单的分析中，我们考虑两个方向相反的主要贡献。当原子彼此接近、相互交叠时，每个电子能够在更大的区域内运动，在其中它的平均势能较低。另一方面，当我们压缩电子气体时，费米能被推高，从而增大了电子的平均动能。在金属密度下，这两种力处于平衡。

%FIG{78}: 内聚的 Wigner–Seitz 模型：(a) 单原子气体；(b) 金属。

为了说明这个论证，设想我们从诸如 Na 这样的碱金属的自由原子出发。每个原子有一个价电子处在原子势 $v_{a}(r)$ 的最低束缚态 $\mathcal{E}_{a}$ 上；在离子实之外，这个势的行为应当如同库仑势 $-e^{2}/r$。固体金属中的传导电子所感受到的晶体势 $\mathcal{V}(\mathbf{r})$ 显然会比这类函数的单纯叠加复杂得多，因为它现在还必须把“其他”传导电子等等的贡献包括进来。然而，作为一种粗糙近似，让我们假定：所选的这个电子凭其静电排斥，暂时把所有其他价电子都从它恰好所处的那个 Wigner–Seitz 原胞（§3.5）的整个范围中排除出去；而晶体中所有其余原胞这时都含有电子，恰好把它们静电中和。换句话说，原子交叠的效应和电子关联等等的效应，都通过取 $\mathcal{V}(\mathbf{r})$ 在每个原胞中等于 $v_{a}(r)$（对 $r\leqslant r_{s}$）来模拟（图 78）。

在原子球的实际边界处，这显然是一个糟糕的近似。尽管如此，其净效果是晶体中电子的平均势能大幅降低。更精确地说，运用 Wigner–Seitz 元胞近似可以证明：导带底 $\mathcal{E}(0)$ 必定位于比这个势垒的能量稍低之处——即位于 $-e^{2}/r_{s}$ 以下的某个能级上。于是当 $r_{s}$ 减小时我们正在赢得势能。

排斥项的理论对自由电子气体来说是不足道的。由 §3.1 的公式我们容易证明：能量态密度（density of states in energy）——单位能量间隔内的电子态数目——由下式给出：
\begin{equation*}
\mathcal{N}(\mathcal{E}_{F})=\frac{3}{2}\frac{n}{\mathcal{E}_{F}}
\tag{4.3}
\end{equation*}
其中 $n$ 是费米能为 $\mathcal{E}_{F}$ 的电子的密度。由此可得
\begin{equation*}
\mathcal{N}(\mathcal{E})\propto\mathcal{E}^{1/2},
\tag{4.4}
\end{equation*}
从而每个电子的平均能量为
\begin{equation*}
\bar{\mathcal{E}}_{\text{kin.}}=\frac{3}{5}\mathcal{E}_{F},
\tag{4.5}
\end{equation*}
它按 $1/r_{s}^{2}$ 变化。

把这两项合并起来，如图 79 那样，我们看到在某个 $r_{s}$ 值处出现一个极小；这个值应当就是实际金属中的原子半径，并且应当给出总内聚能的合理估计。由这个粗糙论证得到的结果至少数量级是对的，说明这个一般性解释在原则上站得住脚。

%FIG{79}: 金属内聚的诸贡献，表为原子间距的函数。

然而有意思的是，所有碱金属的内聚能都非常接近 $\frac{2}{5}\mathcal{E}_{F}$。由这个经验事实，我们从式 (4.5) 推断：金属中的费米能级与能带所由导出的那个原子能级的能量近似重合。这正是 L.C.A.O. 理论中半满能带所应有的结果，如式 (3.30)。对金属内聚的另一种解释则是：电子得以成对地进入 Bloch 态，两电子自旋相反，而不必像它们各自束缚于分立的原子轨道时那样付出静电能量的代价。

内聚力的性质影响固体的弹性和塑性。共价固体刚硬而脆，因为键的方向性阻碍剪切运动，使一个原子无法从另一个原子旁边滑过——正如位错穿过晶格的迁移那样。离子晶体则更具塑性，只要它们是完善的纯晶体（不过普通晶体可能由于生长中带入的缺陷而是脆性的）。静电力没有方向性，因此除了受离子大小的阻碍之外，它们允许离子四处移动。金属是最具塑性的固体，允许位错自由运动，因为电子气体的费米能倾向于使离子彼此保持距离，而内聚能主要是堆积密度的函数。对严格晶格规则性的局域偏离很容易被容纳下来。

\section{刚性能带模型与态密度}

为什么许多金属会以不同的相出现，具有不同的晶体结构，并且随温度变化而从一个相转变到另一个相？与晶体热激发相联系的能量很小——量级为每原子 0·1 eV。不同结构的能量差必定很小。的确，我们上面基于 Wigner–Seitz 公式和紧束缚（tight-binding）公式的粗糙论证，无法区分密度相同而晶体结构不同的相。有意思的是，金属熔化时密度变化不大；上述论证对结构如此不敏感，以至于它们对液态金属同样适用——液态金属可以十分简单地描述为一种无规堆积的球形离子集合体，靠自由电子构成的“胶”黏结在一起。

不过，有一个现象体现了一条关于固体结构能量的原理。考虑下表：

%TAB{}: 代表性合金相的晶体结构与电子/原子比（原书此表未编号）
\begin{table}[H]
\centering
\caption*{合金的晶体结构、代表性合金与电子/原子比}
\begin{tabular}{lccc}
\toprule
结构 & 体心立方 & “$\gamma$ 黄铜” & 六角密堆积 \\
\midrule
合金 & AgZn & Ag$_5$Zn$_8$ & AgZn$_3$ \\
 & Cu$_3$Al & Cu$_9$Al$_4$ & — \\
 & Cu$_5$Sn & Cu$_{31}$Sn$_8$ & Cu$_3$Sn \\
\midrule
电子/原子 & $\frac{3}{2}=1\cdot5$ & $\frac{21}{13}=1\cdot615$ & $\frac{7}{4}=1\cdot75$ \\
\bottomrule
\end{tabular}
\end{table}

每一列指的是一个合金系列，各系列中的合金都具有相同的基本晶体结构，并且各自形成一个成分或多或少确定的稳定相。它们有什么共同之处？在底行我们给出电子数对原子数之比，它对每一列都是常数。例如在 Cu$_9$Al$_4$ 中，Cu 原子提供 9 个电子，4 个 Al 原子每个提供 3 个电子，合计 13 个原子 21 个电子。

这个现象所例示的，就是人们常在合金的刚性能带（rigid band）理论基础上讨论的 Hume–Rothery 定则（Hume–Rothery rules）。其论证是：把组元元素的价电子统统倒入单一的近自由电子（N.F.E.）能带中，这个能带类似于普通纯金属的能带（§3.3）。例如对 Cu$_{31}$Sn$_{8}$，这个假想金属每原子将拥有 1·615 个电子。

我们现在假定合金的布里渊区结构依赖于基本晶格，而不太依赖于占据格点的实际离子。结果表明，对“$\gamma$ 黄铜”结构，这个区的形状和大小恰好使它同一个容纳每原子 1·615 个传导电子的自由电子球相切。于是可以预期，在这些固体中费米面会在相当大的面积上接触区边界，但多半还没有破入下一个区。由于正好在区边界之内的态比具有同一波矢的自由电子能量更低，金属若取这种结构便可在内聚能上获益。关于许多纯金属和复杂合金相之晶体结构的这个解释非常利落，常被算作金属电子理论的一项胜利。遗憾的是，把近自由电子（N.F.E.）模型用于诸如 Cu、Ag、Au 这样的金属是有异议的，因为在纯金属中费米面本来就接触区边界（见 §9.4）。刚性能带模型本身的理论地位也存在一些疑问（见 §5.3），尽管在个别情形中有“区边界效应”的有力证据。

实际上，我们是说：结构使得费米能级落在一大段能量范围的顶端附近，在那里能量态密度很高。这个函数对自由电子气体已由式 (4.3) 给出。要计算实际金属的 $\mathcal{N}(\mathcal{E})$，我们必须求出布里渊区中逐级能量面之间所包围的体积。这个问题在形式上与 §2.5 中晶格谱的计算完全相同。

%FIG{80}: (a) 费米面接触区边界。(b) 能量降到自由电子值之下。(c) 费米能级靠近态密度的极大。

仿照式 (2.66)–(2.69) 的论证，我们可以写出
\begin{equation*}
\mathcal{N}(\mathcal{E})=\frac{1}{4\pi^{3}N\hbar}\int\frac{dS}{|\mathbf{v}_{\mathbf{k}}|},
\tag{4.6}
\end{equation*}
为此我们暂时定义一个具有速度量纲的矢量
\begin{equation*}
\mathbf{v}_{\mathbf{k}}=\frac{1}{\hbar}\nabla_{\mathbf{k}}\mathcal{E}(\mathbf{k}),
\tag{4.7}
\end{equation*}
即能量函数在 $k$ 空间中的梯度。出现在式 (4.6) 中的是这个矢量的模，须在能量为 $\mathcal{E}$ 的能量面上积分。

由于 $\mathcal{E}(\mathbf{k})$ 像 $\nu_{q}$ 一样是倒格子空间中的连续函数，具有倒格子的周期性，van Hove 定理（van Hove's theorem）的论证对它适用。的确，由于 $\mathcal{E}(\mathbf{k})$ 通常在区中心（即 $\nu_{q}$ 中声学支所在的区域）没有特殊的重根，$\mathcal{E}(\mathbf{k})$ 中的奇点必定包括所有四种类型的临界点——极大点、极小点以及两类鞍点。

一般地说，金属的态密度与自由电子值相去不远，只有过渡金属（transition metals）例外：狭窄的 $d$ 带（或几条 $d$ 带）必须能容纳每原子 10 个电子——所以其能量态密度必定远高于普通的 $s$ 带。既有证据表明这一点，也有证据表明 $d$ 电子对金属的内聚有贡献，正如我们从这个模型所预期的。

\section{电子的费米统计}

到目前为止，我们一直假定电子系统处于零温度；按照 Pauli 原理，我们把能级逐一填充，直到把全部电子用完为止，最后填到能量为 $\mathcal{E}_{F}$ 的费米能级。但我们从统计力学知道，电子服从费米–狄拉克统计（费米–狄拉克 statistics）；能量为 $\mathcal{E}$ 的能级被占据的概率是
\begin{equation*}
f^{0}(\mathcal{E})=\frac{1}{e^{(\mathcal{E}-\zeta)/kT}+1},
\tag{4.8}
\end{equation*}
其中 $\zeta$ 是费米势（Fermi potential），即化学势，或每个电子的自由能；只要电子气体处于热力学平衡，它在整个材料中就必须是常数。

设有 $\mathcal{N}(\mathcal{E})\,d\mathcal{E}$ 个态处在能量间隔 $d\mathcal{E}$ 内。若 $n$ 是电子的总密度（每单位晶体体积），则必须有
\begin{equation*}
n=\int_{0}^{\infty}f^{0}(\mathcal{E})\,\mathcal{N}(\mathcal{E})\,d\mathcal{E}.
\tag{4.9}
\end{equation*}
这两个方程定出 $\zeta$ 的值，从而定出整个分布。

在绝对零度下这是不足道的：容易看出，若 $\mathcal{E}<\zeta$，则 $f^{0}(\mathcal{E})$ 等于 1；而当 $\mathcal{E}>\zeta$ 时它为零。于是有
\begin{equation*}
n=\int_{0}^{\zeta}\mathcal{N}(\mathcal{E})\,d\mathcal{E}=\int_{0}^{\mathcal{E}_{F}}\mathcal{N}(\mathcal{E})\,d\mathcal{E},
\tag{4.10}
\end{equation*}
这就是费米能级定义的数学表述：
\begin{equation*}
\zeta=\mathcal{E}_{F}\quad\text{当}\quad T=0.
\tag{4.11}
\end{equation*}

要把 $\zeta$ 表为 $T$ 的函数，原则上需要知道 $\mathcal{N}(\mathcal{E})$。但在金属中分布是高度简并（degenerate）的，即 $kT\ll\mathcal{E}_{F}$。这时我们可以采用如下的数学技巧：——

设有如下形式的积分需要求值：
\begin{equation*}
I=\int_{0}^{\infty}g(\mathcal{E})\,f^{0}(\mathcal{E})\,d\mathcal{E},
\tag{4.12}
\end{equation*}
其中 $g(\mathcal{E})$ 是能量的某个函数。分部积分。我们得到
\begin{equation*}
I=\left[G(\mathcal{E})f^{0}(\mathcal{E})\right]_{0}^{\infty}-\int_{0}^{\infty}G(\mathcal{E})\frac{\partial f^{0}}{\partial\mathcal{E}}\,d\mathcal{E},
\tag{4.13}
\end{equation*}
其中
\begin{equation*}
G(\mathcal{E})\equiv\int_{0}^{\mathcal{E}}g(\mathcal{E})\,d\mathcal{E}.
\tag{4.14}
\end{equation*}
第一项在两个极限下都为零（只要把能量原点取得足够低，就可以令 $g(0)=0$）。

剩下的积分比较容易对付；由于 $\partial f^{0}/\partial\mathcal{E}$ 在 $\mathcal{E}=\zeta$ 处有一个对称的峰，它在那一带的行为宛如一个宽度为 $kT$ 的 $\delta$ 函数（delta-function）。于是，我们把 $G(\mathcal{E})$ 在这一点作 Taylor 展开：
\begin{equation*}
G(\mathcal{E})=G(\zeta)+(\mathcal{E}-\zeta)G'(\zeta)+\tfrac{1}{2}(\mathcal{E}-\zeta)^{2}G''(\zeta)+\dots,
\tag{4.15}
\end{equation*}
并代入式 (4.13)。我们得到
\begin{equation*}
I=G(\zeta)\int_{0}^{\infty}\left(-\frac{\partial f^{0}}{\partial\mathcal{E}}\right)d\mathcal{E}+G'(\zeta)\int_{0}^{\infty}(\mathcal{E}-\zeta)\left(-\frac{\partial f^{0}}{\partial\mathcal{E}}\right)d\mathcal{E}+\dots.
\tag{4.16}
\end{equation*}
但是
\begin{equation*}
\int_{0}^{\infty}\left(-\frac{\partial f^{0}}{\partial\mathcal{E}}\right)d\mathcal{E}=f^{0}(0)-f^{0}(\infty)=1,
\tag{4.17}
\end{equation*}
所以一级近似是
\begin{equation*}
I\approx G(\zeta)=\int_{0}^{\zeta}g(\mathcal{E})\,d\mathcal{E},
\tag{4.18}
\end{equation*}
这正是我们在 $T=0$ 时本应得到的结果。

%FIG{81}: 费米–狄拉克函数及其导数：$T=0$ 时与有限温度下的情形。

式 (4.16) 中的普遍项含有积分
\begin{align*}
F_{n}&=\frac{1}{n!}\int_{0}^{\infty}(\mathcal{E}-\zeta)^{n}\left(-\frac{\partial f^{0}}{\partial\mathcal{E}}\right)d\mathcal{E}\\
&=\frac{(kT)^{n}}{n!}\int_{0}^{\infty}\left(\frac{\mathcal{E}-\zeta}{kT}\right)^{n}\frac{\exp\{(\mathcal{E}-\zeta)/kT\}}{\left[\exp\{(\mathcal{E}-\zeta)/kT\}+1\right]^{2}}\,d\!\left(\frac{\mathcal{E}-\zeta}{kT}\right)\\
&=\frac{(kT)^{n}}{n!}\int_{-\infty}^{\infty}\frac{z^{n}\,dz}{(e^{z}+1)(1+e^{-z})}\\
&=\left\{\begin{array}{ll}
2c_{n}(kT)^{n} & \text{当 } n \text{ 为偶数},\\
0 & \text{当 } n \text{ 为奇数},
\end{array}\right\}
\tag{4.19}
\end{align*}
其中系数 $c_{n}$ 很容易求出，是一些可求和的级数。实践中我们很少超过第二项，
\begin{equation*}
2c_{2}=\frac{1}{6}\pi^{2}.
\tag{4.20}
\end{equation*}
于是
\begin{equation*}
\int_{0}^{\infty}g(\mathcal{E})\,f^{0}(\mathcal{E})\,d\mathcal{E}\approx\int_{0}^{\zeta}g(\mathcal{E})\,d\mathcal{E}+\frac{\pi^{2}}{6}(kT)^{2}\left[\frac{\partial g(\mathcal{E})}{\partial\mathcal{E}}\right]_{\mathcal{E}=\zeta}+\dots.
\tag{4.21}
\end{equation*}
